\section{Introduction}
\label{sec:intro}

Machine agents increasingly carry identity and capability metadata before they
interact~\citep{a2a2026spec}. Regulation already requires providers to identify
AI systems and model versions~\citep{euaiact2024}, while attestation and
provenance infrastructures are being proposed for agent
ecosystems~\citep{chan2024visibility,chan2025infrastructure,shavit2023practices}.
These developments usually frame identity as an accountability tool. Identity
can also change conduct directly: an agent may trust, cooperate with or grant
access to a partner because its credential passes verification.

An identity signal that elicits preferential treatment among its bearers is a
\emph{greenbeard}. Such signals can support cooperation without repeated
interaction, but they fail when the marker separates from the behaviour it is
supposed to identify~\citep{riolo2001evolution,jansen2006altruism,gardner2010greenbeards}.
For machine agents, this separation is an engineering variable rather than a
biological constraint: the pass rate of a forged credential depends on the
verification technology~\citep{jovanovic2024watermarkstealing,anbiaee2026threatmodeling}.
The unresolved question is therefore dynamical. If credentials are costly to
obtain honestly, valuable when accepted and imperfectly verifiable, which
population regimes persist, and how do those regimes change as verification,
penalties and interaction structure vary?

We answer this question by adding a strategic identity layer to a reduced
four-strategy AI race~\citep{domingos2026falling}. Each design combines a badge
that is genuine, forged or absent with conduct toward partners whose badge
passes or fails verification. The construction yields 48 designs evolving on
a 47-dimensional simplex under replicator dynamics; the same payoff structure
also drives a finite-population birth--death process. The underlying race is
held fixed, and every pairwise payoff is evaluated exactly, so all changes are
attributable to the identity layer rather than to sampling or to a modified
interaction game.

The analysis separates local stability from global basin structure. At a
monomorphic resident, invasion is governed by a transversal eigenvalue of the
replicator flow. Most regime boundaries are simple zeros of these eigenvalues
and are therefore available in closed form; numerical integration is used only
to measure basins and to locate attractors. This distinction matters because
the flow is multistable. Averaging states across its disjoint attracting faces
before evaluating a pairwise observable creates cross-regime encounters that
no trajectory contains, so observables must be evaluated within each attractor
and only then averaged over basins.

Three mechanisms organise the results. First, credentials are useful only
where ordinary reciprocity and liability fail to protect safe play. A
reciprocating certified club tolerates more than twice the forgery of an
unconditionally trusting club, showing that in-group conduct is part of the
security design. Penalties help only while detection remains informative: the
protective fine has a vertical asymptote as forged badges become
undetectable. Second, verification does not help a costly club form because an
unbadged population does not read its signal; positive assortment is required
for nucleation. Once formed, however, the club can persist under weaker
conditions than those required to rebuild it after collapse.

Third, certification can fail without an observable deterioration in conduct.
A behavioural mimic copies the club's conduct while avoiding its dues and
crosses its invasion threshold before an exploiter at the baseline. The two
invasion boundaries meet at a codimension-two point, beyond which stronger
assortment removes the visible exploiter while leaving the mimic essentially
unchanged. The credential is then hollow even though behaviour remains safe.
Globally, the dominant trajectories settle on disjoint certified and
uncertified safe faces, while a fully specified targeted start reaches a
numerically robust unsafe mixed rest face outside the uniform sample. Certification
therefore buys robustness to displacement rather than a lower harm level at a
settled state. Its sustaining mechanism is an exclusionary entry barrier: a
club gentle to outsiders is cheaper to imitate and cannot survive forgery.

These results connect nonlinear evolutionary dynamics to certification design.
They also give a practical measurement lesson: conduct alone cannot diagnose
credential integrity. A regime must measure whether passed credentials are
genuine and must treat mutual recognition across providers as part of the
safety architecture, rather than as a separate interoperability choice.
